% =====================================================================
%  L'AIoT & TinyML : quand l'intelligence artificielle s'intègre
%  dans les objets connectés
%  Compilation : latexmk -pdf main.tex   (pdflatex + biber)
% =====================================================================
\documentclass[11pt,a4paper]{report}

% ---------- Langue et encodage ----------
\usepackage[utf8]{inputenc}
\usepackage[T1]{fontenc}
\usepackage[french]{babel}
\usepackage{lmodern}
\usepackage{microtype}
\usepackage[autostyle=true]{csquotes}

% ---------- Mise en page ----------
\usepackage[a4paper,margin=2.2cm,headheight=14pt]{geometry}
\usepackage{setspace}
\setstretch{1.0}
\usepackage{parskip}
\usepackage{fancyhdr}
\usepackage{titlesec}
\usepackage{enumitem}
\setlist{itemsep=2pt,topsep=4pt}

% ---------- Maths, tableaux, figures ----------
\usepackage{amsmath,amssymb}
\usepackage{booktabs}
\usepackage{tabularx}
\usepackage{array}
\usepackage{multirow}
\usepackage{float}
\usepackage{graphicx}
\usepackage[font=small,labelfont=bf]{caption}
\usepackage{xcolor}
\usepackage{tikz}
\usetikzlibrary{arrows.meta,positioning,shapes.geometric,fit,backgrounds,calc}
\usepackage[most]{tcolorbox}

% ---------- Code ----------
\usepackage{listings}

% ---------- Bibliographie ----------
\usepackage[backend=biber,style=numeric,sorting=nyt,defernumbers=true,maxbibnames=4,giveninits=true,url=true,doi=true]{biblatex}
\addbibresource{references.bib}

% ---------- Liens (en dernier) ----------
\usepackage[hidelinks]{hyperref}
\usepackage{bookmark}

% =====================================================================
%  Style
% =====================================================================
\definecolor{bleu}{RGB}{190,79,35}       % orange ULHN
\definecolor{bleuclair}{RGB}{249,234,227}
\definecolor{vert}{RGB}{61,61,60}         % gris anthracite ULHN
\definecolor{vertclair}{RGB}{236,236,235}
\definecolor{orange}{RGB}{140,58,26}      % orange ULHN foncé
\definecolor{orangeclair}{RGB}{243,224,214}
\definecolor{gris}{RGB}{110,110,110}
\definecolor{fondcode}{RGB}{246,246,246}

% Titres de chapitres compacts
\titleformat{\chapter}[display]
  {\normalfont\bfseries\color{bleu}}
  {\large\scshape Partie~\thechapter}
  {4pt}
  {\LARGE}[\vspace{2pt}{\color{bleu}\titlerule[1pt]}]
\titlespacing*{\chapter}{0pt}{-20pt}{18pt}
\titleformat{\section}{\normalfont\Large\bfseries\color{bleu}}{\thesection}{0.8em}{}
\titleformat{\subsection}{\normalfont\large\bfseries\color{bleu!80!black}}{\thesubsection}{0.8em}{}
\titleformat{\subsubsection}{\normalfont\normalsize\bfseries}{\thesubsubsection}{0.8em}{}
\setcounter{secnumdepth}{3}
\setcounter{tocdepth}{2}

% En-têtes et pieds de page
\pagestyle{fancy}
\fancyhf{}
\fancyhead[L]{\small\color{gris}AIoT \& TinyML}
\fancyhead[R]{\small\color{gris}\nouppercase{\leftmark}}
\fancyfoot[C]{\thepage}
\renewcommand{\headrulewidth}{0.4pt}
\renewcommand{\chaptermark}[1]{\markboth{Partie~\thechapter.~#1}{}}
\fancypagestyle{plain}{\fancyhf{}\fancyfoot[C]{\thepage}\renewcommand{\headrulewidth}{0pt}}

% Listings
\lstdefinestyle{code}{
  backgroundcolor=\color{fondcode},
  basicstyle=\ttfamily\small,
  keywordstyle=\color{bleu}\bfseries,
  commentstyle=\color{vert}\itshape,
  stringstyle=\color{orange},
  frame=single, rulecolor=\color{gray!50},
  breaklines=true, showstringspaces=false,
  numbers=left, numberstyle=\tiny\color{gris}, xleftmargin=1.5em,
  columns=fullflexible, keepspaces=true,
  literate={é}{{\'e}}1 {è}{{\`e}}1 {à}{{\`a}}1 {ê}{{\^e}}1 {ç}{{\c c}}1 {ô}{{\^o}}1 {î}{{\^i}}1 {É}{{\'E}}1
}
\lstset{style=code}
\renewcommand{\lstlistingname}{Listing}

% Encadrés
\newtcolorbox{encadre}[1][]{colback=bleuclair,colframe=bleu,boxrule=0.6pt,arc=2pt,
  left=6pt,right=6pt,top=4pt,bottom=4pt,fonttitle=\bfseries,#1}
\newtcolorbox{aretenir}{colback=vertclair,colframe=vert,boxrule=0.6pt,arc=2pt,
  left=6pt,right=6pt,top=4pt,bottom=4pt,title={À retenir},fonttitle=\bfseries}

% Styles TikZ communs
\tikzset{
  bloc/.style={draw=bleu, fill=bleuclair, rounded corners=3pt, align=center,
               minimum height=1cm, inner sep=4pt, font=\small},
  blocv/.style={draw=vert, fill=vertclair, rounded corners=3pt, align=center,
               minimum height=1cm, inner sep=4pt, font=\small},
  bloco/.style={draw=orange, fill=orangeclair, rounded corners=3pt, align=center,
               minimum height=1cm, inner sep=4pt, font=\small},
  fleche/.style={-{Stealth[length=2.2mm]}, thick, draw=gris},
}

\newcommand{\code}[1]{\texttt{#1}}

% =====================================================================
\begin{document}

% ---------------- Page de titre ----------------
\begin{titlepage}
  \centering
  \includegraphics[width=5.5cm]{images/logo_ulhn.png}\par
  \vspace{1.2cm}
  {\large\scshape Module : Internet des Objets\par}
  \vspace{0.4cm}
  {\color{bleu}\rule{\linewidth}{1.2pt}\par}
  \vspace{0.6cm}
  {\Huge\bfseries\color{bleu} L'AIoT \& le TinyML\par}
  \vspace{0.4cm}
  {\LARGE Quand l'intelligence artificielle\\ s'intègre dans les objets connectés\par}
  \vspace{0.6cm}
  {\color{bleu}\rule{\linewidth}{1.2pt}\par}
  \vspace{1.5cm}
  {\Large Rapport de veille documentaire\par}
  \vspace{2cm}
  \begin{minipage}[t]{0.48\linewidth}
    \raggedright
    {\large\bfseries Réalisé par :\par}
    \vspace{0.2cm}
    {\large Hamadou BA\\ Mohamed SANOGO\\ Nicolas CHEVALIER\par}
  \end{minipage}\hfill
  \begin{minipage}[t]{0.48\linewidth}
    \raggedleft
    {\large\bfseries Sous la direction de :\par}
    \vspace{0.2cm}
    {\large M. Claude DUVALLET\par}
  \end{minipage}
  \vfill
  {\large Université Le Havre Normandie\par}
  \vspace{0.2cm}
  {\large Année universitaire 2026--2027\par}
\end{titlepage}

% ---------------- Résumé ----------------
\pagenumbering{roman}
\chapter*{Résumé}
\addcontentsline{toc}{chapter}{Résumé}
\markboth{Résumé}{}

L'Internet des objets (IoT) a permis de connecter des milliards de capteurs et
d'actionneurs au réseau. Pendant longtemps, ces objets se sont contentés de
mesurer puis de transmettre : l'analyse des données et la prise de décision
étaient confiées à des serveurs distants, dans le cloud. L'\emph{Artificial
Intelligence of Things} (AIoT) désigne la convergence de l'IoT et de
l'intelligence artificielle, et le \emph{Tiny Machine Learning} (TinyML) en
constitue la forme la plus contrainte : l'exécution de modèles d'apprentissage
automatique directement sur des microcontrôleurs disposant de quelques centaines
de kilooctets de mémoire et consommant quelques milliwatts.

Ce rapport présente d'abord les notions d'IoT, d'AIoT et de TinyML, ainsi que
l'architecture qui les relie. Il détaille ensuite les principes techniques qui
rendent le TinyML possible (pipeline de développement, quantification, élagage,
distillation), le matériel et les outils logiciels disponibles, puis les
bénéfices, les domaines d'application, les limites et les perspectives de cette
approche. Le fil conducteur est un cas d'usage concret : la détection d'anomalies
de vibration sur un moteur à l'aide d'un ESP32.

\textbf{Mots-clés :} Internet des objets, AIoT, TinyML, Edge AI, microcontrôleur,
ESP32, quantification, inférence embarquée.

\tableofcontents

% ---------------- Introduction générale ----------------
\clearpage
\pagenumbering{arabic}
% =====================================================================
%  Introduction générale
% =====================================================================
\chapter*{Introduction générale}
\addcontentsline{toc}{chapter}{Introduction générale}
\markboth{Introduction générale}{}

En une quinzaine d'années, l'Internet des objets (\emph{Internet of Things}, IoT)
est passé du statut de concept à celui d'infrastructure du quotidien :
thermostats, montres, compteurs électriques, machines industrielles ou capteurs
agricoles sont désormais connectés. Le cabinet IoT Analytics estimait à environ
18,8~milliards le nombre d'objets connectés actifs dans le monde fin
2024~\cite{iotanalytics2024}. Ces objets produisent un volume de données
considérable, que l'architecture « classique » de l'IoT envoie vers le cloud pour
y être stockées, analysées et transformées en décisions.

Ce modèle centralisé montre aujourd'hui ses limites. Transmettre en permanence
des mesures brutes consomme de la bande passante et de l'énergie, introduit une
latence qui dépend de la qualité du réseau, rend le système inopérant lorsque la
connexion est coupée et expose des données parfois sensibles (sons, images,
mesures physiologiques). Le paradigme de l'\emph{edge computing} propose de
rapprocher le traitement de la source des données~\cite{shi2016edge}.

L'étape suivante consiste à doter les objets eux-mêmes d'une capacité
d'analyse : c'est l'\textbf{AIoT}, contraction d'\emph{Artificial Intelligence}
et d'\emph{Internet of Things}~\cite{zhang2021aiot}. Lorsque cette intelligence
est logée dans les composants les plus modestes de la chaîne, des
microcontrôleurs de quelques euros fonctionnant sur pile, on parle de
\textbf{TinyML}~\cite{warden2019tinyml,banbury2020benchmarking}. Un capteur
de vibrations ne transmet plus des milliers d'échantillons par seconde : il
signale simplement « fonctionnement anormal détecté ».

\begin{encadre}[title={Problématique}]
Comment faire fonctionner des modèles d'intelligence artificielle, conçus
habituellement pour des serveurs dotés de gigaoctets de mémoire, sur des objets
connectés qui ne disposent que de quelques centaines de kilooctets, et quels
bénéfices et limites cette approche apporte-t-elle aux systèmes IoT ?
\end{encadre}

Pour y répondre, ce rapport est organisé en cinq parties. La
\textbf{partie~1} pose les bases : objet connecté, architecture d'un système IoT,
définition de l'AIoT et du TinyML, et comparaison entre une architecture centrée
sur le cloud et une architecture où l'inférence est réalisée localement. La
\textbf{partie~2} détaille les principes techniques du TinyML, en particulier
le pipeline de développement et les techniques de compression des modèles. La
\textbf{partie~3} présente le matériel et la chaîne d'outils logiciels. La
\textbf{partie~4} examine les bénéfices de l'approche et ses domaines
d'application. Enfin, la \textbf{partie~5} discute les limites, les défis et les
perspectives d'avenir.


% ---------------- Parties ----------------
% =====================================================================
%  PARTIE 1 — De l'IoT à l'AIoT et au TinyML
% =====================================================================
\chapter{De l'IoT à l'AIoT et au TinyML}
\label{ch:fondements}

Avant d'étudier la manière dont un modèle d'intelligence artificielle peut être
exécuté sur un microcontrôleur, il est nécessaire de préciser le vocabulaire et
le cadre dans lequel s'inscrit ce travail. Cette partie définit successivement
l'objet connecté et l'architecture d'un système IoT
(section~\ref{sec:iot}), l'AIoT (section~\ref{sec:aiot}), le TinyML
(section~\ref{sec:tinyml}), puis compare les architectures centrées sur le cloud
à celles où l'inférence est réalisée dans l'objet
(section~\ref{sec:archi}).

% ---------------------------------------------------------------------
\section{L'Internet des objets}
\label{sec:iot}

\subsection{Qu'est-ce qu'un objet connecté ?}

L'expression \emph{Internet of Things} est attribuée à Kevin Ashton, qui
l'utilise dès 1999 pour décrire des objets capables de collecter eux-mêmes des
informations sur le monde physique, sans saisie humaine~\cite{ashton2009}. La
recommandation ITU-T~Y.2060 de l'Union internationale des télécommunications en
donne une définition de référence : l'IoT est \og une infrastructure mondiale
pour la société de l'information, qui permet de disposer de services évolués en
interconnectant des objets (physiques ou virtuels) grâce aux technologies de
l'information et de la communication interopérables existantes ou en
évolution\fg{}~\cite{itu2012y2060}.

Un \textbf{objet connecté} est donc un dispositif physique qui associe :
\begin{itemize}
  \item un ou plusieurs \textbf{capteurs} qui mesurent une grandeur physique
        (température, accélération, son, lumière, humidité, pression\dots) ;
  \item une \textbf{unité de traitement}, le plus souvent un microcontrôleur
        (MCU), qui pilote les capteurs et met en forme les données ;
  \item un \textbf{module de communication} (Wi-Fi, Bluetooth Low Energy,
        LoRaWAN, NB-IoT, Zigbee\dots) qui relie l'objet au réseau ;
  \item éventuellement des \textbf{actionneurs} (relais, moteur, LED,
        buzzer) qui agissent sur l'environnement ;
  \item une \textbf{source d'énergie}, souvent une batterie ou une pile, parfois
        complétée par de la récupération d'énergie (panneau solaire, vibrations).
\end{itemize}

La carte ESP32 d'Espressif illustre bien ce type de plateforme : un
microcontrôleur cadencé à 240~MHz, environ 520~Ko de SRAM, du Wi-Fi et du
Bluetooth intégrés, pour un coût de quelques euros~\cite{espressif-esp32}.
La contrainte majeure de ces objets n'est pas tant leur puissance de calcul que
leur \textbf{budget énergétique} : un capteur installé dans un champ ou sur une
machine doit souvent fonctionner plusieurs mois, voire plusieurs années, sans
intervention.

\subsection{Architecture générale d'un système IoT}

Il n'existe pas d'architecture IoT unique, mais la littérature converge vers un
découpage en couches~\cite{alfuqaha2015,atzori2010}. Le modèle le plus courant
distingue quatre niveaux, représentés à la figure~\ref{fig:archi-iot} :

\begin{enumerate}
  \item \textbf{La couche perception} (ou couche objets) regroupe les capteurs et
        actionneurs. Elle convertit des phénomènes physiques en signaux
        numériques et, inversement, exécute des commandes.
  \item \textbf{La couche réseau} (ou transport) achemine les données. Elle
        s'appuie sur des protocoles radio adaptés aux faibles débits et à la
        faible consommation (BLE, Zigbee, LoRaWAN, NB-IoT) ainsi que sur des
        passerelles (\emph{gateways}) qui relaient les messages vers Internet.
  \item \textbf{La couche traitement} (ou plateforme, \emph{middleware}) stocke,
        agrège et analyse les données, le plus souvent dans le cloud. On y trouve
        les bases de données, les tableaux de bord et, aujourd'hui, les modèles
        d'apprentissage automatique.
  \item \textbf{La couche application} fournit le service final à l'utilisateur :
        application mobile, alerte de maintenance, pilotage d'un système
        d'irrigation, etc.
\end{enumerate}

\begin{figure}[H]
\centering
\begin{tikzpicture}[
    couche/.style={draw=#1, fill=#1!12, rounded corners=3pt, minimum width=13.6cm,
                   minimum height=1.25cm},
    lbl/.style={font=\bfseries\small, anchor=west, text width=3cm},
    ex/.style={font=\small, anchor=west, text width=9.6cm}]
  \node[couche=orange] (app) at (0,4.5) {};
  \node[lbl] at ($(app.west)+(0.2,0)$) {4. Application};
  \node[ex]  at ($(app.west)+(3.6,0)$) {Tableaux de bord, applications mobiles, alertes,
       maintenance, villes et maisons intelligentes};
  \node[couche=vert] (trt) at (0,3) {};
  \node[lbl] at ($(trt.west)+(0.2,0)$) {3. Traitement};
  \node[ex]  at ($(trt.west)+(3.6,0)$) {Cloud / serveurs : stockage, base de données,
       analyse, entraînement des modèles d'IA};
  \node[couche=bleu] (res) at (0,1.5) {};
  \node[lbl] at ($(res.west)+(0.2,0)$) {2. Réseau};
  \node[ex]  at ($(res.west)+(3.6,0)$) {Wi-Fi, BLE, Zigbee, LoRaWAN, NB-IoT, passerelles ;
       protocoles MQTT, CoAP, HTTP};
  \node[couche=gray] (per) at (0,0) {};
  \node[lbl] at ($(per.west)+(0.2,0)$) {1. Perception};
  \node[ex]  at ($(per.west)+(3.6,0)$) {Capteurs (température, accéléromètre, micro\dots),
       actionneurs, microcontrôleurs};
  \draw[fleche] ($(per.east)+(0.3,0)$) -- node[right,font=\footnotesize,align=left]{données\\montantes} ($(app.east)+(0.3,0)$);
\end{tikzpicture}
\caption{Architecture en quatre couches d'un système IoT (d'après~\cite{alfuqaha2015}).}
\label{fig:archi-iot}
\end{figure}

Dans ce schéma classique, l'« intelligence » du système réside dans la couche
traitement. L'objet est un simple producteur de données : il mesure, transmet et
attend éventuellement un ordre en retour. Cette organisation a l'avantage de la
simplicité et de la puissance (le cloud dispose de ressources pratiquement
illimitées), mais elle fait dépendre toute décision de la disponibilité et de la
qualité du réseau.

\subsection{Les limites du modèle « tout cloud »}

Plusieurs difficultés apparaissent lorsque le nombre d'objets et le volume de
données augmentent~\cite{shi2016edge} :
\begin{itemize}
  \item \textbf{La bande passante} : un accéléromètre échantillonné à 1~kHz sur
        trois axes en 16~bits produit environ 6~Ko/s, soit plus de 500~Mo par
        jour et par capteur. Multiplier ce flux par des centaines de machines
        sature rapidement un réseau LoRaWAN ou NB-IoT, dont les débits se
        comptent en kilobits par seconde.
  \item \textbf{La latence} : l'aller-retour vers un serveur ajoute des délais
        variables, parfois de plusieurs centaines de millisecondes, incompatibles
        avec certaines réactions (arrêt d'une machine, détection d'une chute).
  \item \textbf{L'énergie} : sur un objet sur batterie, la transmission radio
        est souvent le poste de consommation le plus important, bien devant le
        calcul~\cite{warden2019tinyml}.
  \item \textbf{La confidentialité} : envoyer en continu le son d'une pièce ou
        le rythme cardiaque d'une personne multiplie les risques de fuite.
  \item \textbf{La disponibilité} : sans connexion, l'objet devient « aveugle ».
\end{itemize}
Ces limites motivent le déplacement d'une partie de l'analyse vers la
périphérie du réseau (\emph{edge}), puis vers l'objet lui-même.

% ---------------------------------------------------------------------
\section{L'AIoT : l'intelligence artificielle au service des objets}
\label{sec:aiot}

\subsection{Définition}

L'\textbf{AIoT} (\emph{Artificial Intelligence of Things}) désigne la
combinaison des technologies de l'intelligence artificielle, en particulier de
l'apprentissage automatique, avec l'infrastructure de l'Internet des
objets~\cite{zhang2021aiot}. L'objectif est de transformer les données
collectées par les objets en informations exploitables (reconnaître une
situation, détecter une anomalie, anticiper un événement) et, idéalement, de
déclencher une action sans intervention humaine.

L'IA peut être placée à différents endroits de l'architecture : dans le
cloud, sur une passerelle ou un serveur de proximité (\emph{edge server}), ou
directement dans l'objet (\emph{on-device AI}). L'AIoT recouvre l'ensemble de
ces configurations ; ce qui la caractérise est le passage d'objets qui
\emph{mesurent} à des systèmes qui \emph{comprennent} et \emph{décident}.

\subsection{Pourquoi ajouter de l'IA aux objets connectés ?}

Les règles fixes (\og si la température dépasse 80~°C, alors alarme\fg{})
suffisent pour des phénomènes simples. Elles deviennent inadaptées lorsque le
signal est complexe : reconnaître un mot prononcé, distinguer le bruit normal
d'un roulement de celui d'un roulement usé, identifier un geste à partir d'un
accéléromètre. Dans ces cas, il est très difficile d'écrire les règles à la
main, alors qu'un modèle d'apprentissage automatique peut les apprendre à
partir d'exemples. L'IA apporte ainsi aux systèmes IoT :
\begin{itemize}
  \item la \textbf{reconnaissance de motifs} dans des signaux bruités (audio,
        vibrations, images) ;
  \item la \textbf{détection d'anomalies}, c'est-à-dire l'identification
        d'écarts par rapport à un comportement appris comme normal ;
  \item la \textbf{prédiction} de l'évolution d'une grandeur (consommation,
        usure, rendement) ;
  \item l'\textbf{adaptation} au contexte et à l'utilisateur.
\end{itemize}

\subsection{Exemples d'utilisation}

\begin{table}[H]
\centering
\small
\begin{tabularx}{\linewidth}{>{\bfseries}l X X}
\toprule
Domaine & Données collectées & Apport de l'IA \\
\midrule
Industrie      & Vibrations, courant, température des moteurs & Détection d'anomalies, maintenance prédictive \\
Maison         & Son, présence, consommation électrique & Reconnaissance de mots-clés, automatisation \\
Agriculture    & Humidité du sol, météo, images de feuilles & Irrigation optimisée, détection de maladies \\
Santé          & Fréquence cardiaque, mouvements & Détection de chutes ou d'arythmies \\
Transport      & Accélérations, images, données moteur & Aide à la conduite, surveillance du véhicule \\
Environnement  & Sons de la forêt, qualité de l'air & Détection de tronçonneuses, alerte pollution \\
\bottomrule
\end{tabularx}
\caption{Exemples d'applications de l'AIoT.}
\label{tab:exemples-aiot}
\end{table}

Ces exemples, développés dans la partie~\ref{ch:benefices}, ont un point
commun : la décision utile (\og anomalie\fg{}, \og mot-clé reconnu\fg{},
\og chute\fg{}) représente une quantité d'information infime par rapport aux
données brutes dont elle est issue. C'est précisément ce qui rend intéressante
l'exécution du modèle au plus près du capteur.

% ---------------------------------------------------------------------
\newpage
\section{Le TinyML : l'apprentissage automatique sur microcontrôleur}
\label{sec:tinyml}

\subsection{Définition}

Le \textbf{TinyML} (\emph{Tiny Machine Learning}) désigne le domaine qui
regroupe les technologies, matérielles et logicielles, permettant d'exécuter
des modèles d'apprentissage automatique sur des appareils à très faibles
ressources, typiquement des microcontrôleurs, avec une consommation de l'ordre
du milliwatt ou moins~\cite{warden2019tinyml,banbury2020benchmarking}. Le
terme a été popularisé en 2019 par l'ouvrage de Pete Warden et Daniel
Situnayake~\cite{warden2019tinyml} et par la création de la communauté TinyML
Foundation, devenue depuis EDGE~AI FOUNDATION~\cite{tinymlfoundation}.

Le TinyML se situe à l'extrémité du continuum de l'IA embarquée, comme le
montre le tableau~\ref{tab:spectre}. Là où un serveur dispose de dizaines de
gigaoctets de mémoire et un smartphone de plusieurs gigaoctets, un
microcontrôleur n'en possède que quelques centaines de kilooctets, soit un
facteur de l'ordre de $10^5$.

\begin{table}[H]
\centering
\small
\begin{tabularx}{\linewidth}{l X X X}
\toprule
 & \textbf{Cloud AI} & \textbf{Mobile / Edge AI} & \textbf{TinyML} \\
\midrule
Matériel        & Serveurs GPU/TPU & Smartphone, Raspberry Pi, Jetson & Microcontrôleur (ESP32, Arm Cortex-M) \\
Mémoire vive    & Dizaines à centaines de Go & 1 à 16 Go & 32 Ko à quelques Mo \\
Stockage        & To & Dizaines de Go & 256 Ko à quelques Mo de flash \\
Puissance       & Centaines de watts & Quelques watts & Milliwatts \\
Système         & Linux, conteneurs & Android, Linux & Pas d'OS ou RTOS (FreeRTOS) \\
Coût unitaire   & Milliers d'euros & 50 à 500 € & 1 à 10 € \\
Taille de modèle & Go & Dizaines de Mo & Ko à quelques centaines de Ko \\
\bottomrule
\end{tabularx}
\caption{Positionnement du TinyML par rapport aux autres formes d'IA (ordres de grandeur, d'après~\cite{banbury2020benchmarking,lin2020mcunet}).}
\label{tab:spectre}
\end{table}

\subsection{Les contraintes du TinyML}

Concevoir une application TinyML revient à respecter simultanément plusieurs
contraintes très fortes~\cite{banbury2020benchmarking,ray2022review} :

\paragraph{La mémoire.} Un microcontrôleur dispose de deux types de mémoire.
La \textbf{mémoire flash} (non volatile) stocke le programme et les poids du
modèle ; la \textbf{SRAM} (volatile) contient les données de travail, en
particulier les activations intermédiaires calculées pendant l'inférence. Les
deux sont limitées : un modèle doit typiquement tenir dans quelques centaines de
kilooctets de flash et ses activations dans quelques dizaines à centaines de
kilooctets de SRAM. Lin \emph{et al.} montrent que c'est souvent la SRAM, plus
encore que la flash, qui constitue le goulot d'étranglement~\cite{lin2020mcunet}.

\paragraph{La puissance de calcul.} Les microcontrôleurs fonctionnent à des
fréquences de quelques dizaines à quelques centaines de mégahertz, souvent sans
unité de calcul flottant performante. Un réseau de neurones de quelques
millions d'opérations est envisageable ; un modèle de plusieurs milliards
d'opérations ne l'est pas.

\paragraph{L'énergie.} Pour fonctionner des mois sur une pile, l'objet doit
passer l'essentiel de son temps en veille et limiter la durée de chaque
inférence. La consommation se mesure alors en milliwatts et l'énergie par
inférence en microjoules ou millijoules~\cite{banbury2021mlperftiny}.

\paragraph{L'environnement logiciel.} Pas de système d'exploitation complet,
pas toujours d'allocation dynamique de mémoire, pas de Python : le modèle est
compilé avec le micrologiciel (\emph{firmware}) en C/C++ et exécuté par un
moteur d'inférence minimaliste comme TensorFlow Lite for Microcontrollers
(TFLM)~\cite{david2021tflm}.

\subsection{Entraînement et inférence : une séparation fondamentale}

Il est essentiel de distinguer les deux phases de vie d'un modèle :
\begin{itemize}
  \item l'\textbf{entraînement}, qui consiste à ajuster les paramètres du modèle
        à partir d'un grand nombre d'exemples. Il demande beaucoup de mémoire
        et de calcul et s'effectue sur un ordinateur ou dans le cloud ;
  \item l'\textbf{inférence}, qui consiste à appliquer le modèle déjà entraîné à
        une nouvelle donnée pour obtenir une prédiction. C'est la seule phase
        exécutée sur le microcontrôleur dans la grande majorité des systèmes
        TinyML~\cite{david2021tflm}.
\end{itemize}
Entre les deux, le modèle doit être \textbf{optimisé} (réduit, quantifié) puis
\textbf{converti} dans un format exécutable par le microcontrôleur. Ces étapes
sont détaillées dans la partie~\ref{ch:principes}.

\begin{aretenir}
L'IoT connecte les objets, l'AIoT leur donne la capacité d'interpréter leurs
données, et le TinyML pousse cette logique jusqu'aux plus petits composants :
des microcontrôleurs à quelques euros, alimentés sur pile, qui exécutent
localement un modèle entraîné ailleurs.
\end{aretenir}

% ---------------------------------------------------------------------
\section{Fonctionnement d'une architecture AIoT/TinyML}
\label{sec:archi}

\subsection{La chaîne capteur -- microcontrôleur -- modèle -- action}

Dans un système TinyML, le traitement suit une boucle locale composée de cinq
étapes (figure~\ref{fig:boucle}) :

\begin{enumerate}
  \item \textbf{Acquisition} : le capteur mesure le phénomène (par exemple
        l'accélération d'un moteur) et le microcontrôleur lit les échantillons
        via un bus I\textsuperscript{2}C, SPI ou un convertisseur
        analogique-numérique.
  \item \textbf{Prétraitement} : les échantillons sont regroupés dans une
        fenêtre temporelle (par exemple une seconde) puis transformés en
        caractéristiques plus compactes : moyenne, écart-type, valeur efficace,
        spectre par transformée de Fourier (FFT), coefficients MFCC pour l'audio.
  \item \textbf{Inférence} : le vecteur de caractéristiques est fourni au modèle,
        qui renvoie une prédiction (classe, score d'anomalie, valeur).
  \item \textbf{Décision} : la prédiction est comparée à un seuil ou interprétée
        par une règle simple.
  \item \textbf{Action et communication} : l'objet agit (LED, relais, arrêt
        d'urgence) et ne transmet au réseau que l'information utile, par exemple
        un message MQTT \og anomalie détectée\fg{}.
\end{enumerate}

\begin{figure}[H]
\centering
\begin{tikzpicture}[node distance=0.55cm]
  \node[bloc, minimum width=2.2cm] (cap) {Capteur\\\footnotesize(accéléromètre)};
  \node[bloc, minimum width=2.2cm, right=of cap] (pre) {Prétraitement\\\footnotesize(fenêtre, FFT)};
  \node[blocv, minimum width=2.2cm, right=of pre] (inf) {Inférence\\\footnotesize(modèle TinyML)};
  \node[blocv, minimum width=2.0cm, right=of inf] (dec) {Décision\\\footnotesize(seuil)};
  \node[bloco, minimum width=2.0cm, right=of dec] (act) {Action\\\footnotesize(alerte, LED)};
  \draw[fleche] (cap) -- (pre);
  \draw[fleche] (pre) -- (inf);
  \draw[fleche] (inf) -- (dec);
  \draw[fleche] (dec) -- (act);
  \begin{scope}[on background layer]
    \node[draw=gris, dashed, rounded corners=6pt, fit=(pre)(inf)(dec), inner sep=8pt,
          label={[font=\footnotesize\itshape,text=gris]below:dans le microcontrôleur (ESP32)}] {};
  \end{scope}
  \node[bloco, minimum width=2.6cm, above=1.1cm of act, xshift=-1.2cm] (cloud) {Cloud / tableau de bord\\\footnotesize(alertes uniquement)};
  \draw[fleche, dashed] (act) -- node[right,font=\footnotesize]{MQTT} (cloud);
\end{tikzpicture}
\caption{Boucle de traitement locale d'un objet AIoT/TinyML.}
\label{fig:boucle}
\end{figure}

Le listing~\ref{lst:boucle} montre, sous forme simplifiée, la structure du
programme principal d'un ESP32 qui applique ce principe. Le modèle est
interrogé à chaque fenêtre de mesures et seule une alerte est publiée.

\begin{lstlisting}[language=C++, caption={Structure simplifiée de la boucle d'inférence sur ESP32.}, label={lst:boucle}]
void loop() {
  // 1. Acquisition : une fenêtre de 1 s à 1 kHz sur 3 axes
  lire_accelerometre(fenetre, 1000);

  // 2. Prétraitement : extraction de caractéristiques (RMS, FFT...)
  calculer_caracteristiques(fenetre, entree_modele);

  // 3. Inférence locale (TensorFlow Lite for Microcontrollers)
  interpreter->Invoke();
  float score = sortie_modele[0];

  // 4-5. Décision et action : on ne transmet que l'événement utile
  if (score > SEUIL_ANOMALIE) {
    digitalWrite(LED_ALERTE, HIGH);
    mqtt.publish("usine/moteur1/alerte", "anomalie");
  }
  esp_light_sleep_start();   // mise en veille jusqu'à la prochaine fenêtre
}
\end{lstlisting}

\subsection{Comparaison avec une architecture centrée sur le cloud}

Dans une architecture IoT classique, l'objet transmet en continu ses mesures
brutes ; la passerelle les relaie vers le cloud, où le modèle est exécuté avant
que la décision ne redescende éventuellement vers l'objet. Dans une architecture
TinyML, l'inférence est réalisée dans l'objet et seules les décisions remontent.
La figure~\ref{fig:cloud-vs-edge} compare les deux flux.

\begin{figure}[H]
\centering
\begin{tikzpicture}[node distance=0.9cm]
  % ---- Architecture cloud
  \node[font=\bfseries\small] (t1) at (0,0) {a) Architecture centrée sur le cloud};
  \node[bloc, below=0.4cm of t1, xshift=-5cm, minimum width=2cm] (c1) {Capteur\\+ MCU};
  \node[bloc, right=of c1, minimum width=2cm] (g1) {Passerelle};
  \node[bloco, right=of g1, minimum width=2.4cm] (cl1) {Cloud\\(modèle IA)};
  \node[bloc, right=of cl1, minimum width=2cm] (a1) {Décision\\/ action};
  \draw[fleche, draw=orange, line width=1.6pt] (c1) -- node[above,font=\scriptsize]{brut} (g1);
  \draw[fleche, draw=orange, line width=1.6pt] (g1) -- node[above,font=\scriptsize]{brut} (cl1);
  \draw[fleche] (cl1) -- (a1);
  \draw[fleche, dashed] (a1.south) -- ++(0,-0.45) -| node[pos=0.25,below,font=\scriptsize]{commande en retour (latence réseau)} (c1.south);

  % ---- Architecture TinyML
  \node[font=\bfseries\small, below=2.4cm of t1] (t2) {b) Architecture TinyML (inférence locale)};
  \node[blocv, below=0.4cm of t2, xshift=-5cm, minimum width=2cm] (c2) {Capteur\\+ MCU\\+ modèle};
  \node[blocv, right=of c2, minimum width=2cm] (a2) {Décision\\/ action\\locale};
  \node[bloc, right=of a2, minimum width=2cm] (g2) {Passerelle};
  \node[bloc, right=of g2, minimum width=2.4cm] (cl2) {Cloud\\(supervision)};
  \draw[fleche] (c2) -- (a2);
  \draw[fleche, draw=vert, line width=0.6pt] (a2) -- node[above,font=\scriptsize]{alerte} (g2);
  \draw[fleche, draw=vert, line width=0.6pt] (g2) -- node[above,font=\scriptsize]{alerte} (cl2);
\end{tikzpicture}
\caption{Flux de données : architecture cloud (données brutes, flèches épaisses) et architecture TinyML (seules les alertes remontent).}
\label{fig:cloud-vs-edge}
\end{figure}

\begin{table}[H]
\centering
\small
\begin{tabularx}{\linewidth}{>{\bfseries}l X X}
\toprule
Critère & Architecture cloud & Architecture TinyML \\
\midrule
Lieu de l'inférence      & Serveur distant & Microcontrôleur de l'objet \\
Données transmises       & Mesures brutes, en continu & Résultats et alertes uniquement \\
Latence de décision      & Dépend du réseau (de dizaines de ms à plusieurs secondes) & Quelques ms à quelques dizaines de ms, indépendante du réseau \\
Fonctionnement hors ligne & Impossible & Possible pour les fonctions locales \\
Confidentialité          & Données sensibles transférées et stockées à distance & Données brutes conservées dans l'objet \\
Complexité des modèles   & Pratiquement illimitée & Fortement contrainte (Ko de mémoire) \\
Mise à jour du modèle    & Simple (côté serveur) & Plus délicate (mise à jour du firmware, OTA) \\
Coût récurrent           & Abonnement cloud, transfert de données & Faible, surtout lié à la supervision \\
\bottomrule
\end{tabularx}
\caption{Comparaison entre une architecture centrée sur le cloud et une architecture TinyML.}
\label{tab:cloud-vs-tinyml}
\end{table}

Les deux approches ne s'opposent pas nécessairement. Dans la pratique, les
architectures AIoT sont souvent \textbf{hybrides} : l'objet filtre et
interprète localement les données, une passerelle ou le cloud agrège les
résultats de nombreux objets, réalise les analyses lourdes et ré-entraîne
périodiquement les modèles, qui sont ensuite redéployés dans les objets. Cette
complémentarité est discutée dans la partie~\ref{ch:limites}.

\subsection{Le cas d'étude retenu}

Pour illustrer concrètement ces principes, ce rapport s'appuie sur un exemple
fil rouge : la \textbf{détection d'anomalies vibratoires sur un moteur}. Un
accéléromètre (de type MPU6050 ou ADXL345) est fixé sur le carter du moteur et
relié à un ESP32. Un petit modèle, entraîné sur ordinateur à partir de mesures
en fonctionnement normal, est quantifié puis intégré au firmware. L'ESP32 analyse
chaque seconde de vibrations et publie une alerte lorsque le comportement
s'écarte de la normale. Ce scénario mobilise toutes les briques étudiées dans
les parties suivantes : pipeline de développement, compression du modèle,
matériel et outils, bénéfices et limites.

% =====================================================================
%  PARTIE 2 — Principes techniques du TinyML
% =====================================================================
\chapter{Principes techniques du TinyML}
\label{ch:principes}

\section{Introduction}

La partie précédente a permis de définir le concept de TinyML et d'en
identifier les enjeux fondamentaux. Cette partie examine les mécanismes
techniques sous-jacents rendant possible l'exécution de modèles d'intelligence
artificielle sur des dispositifs embarqués aux ressources limitées.

La problématique centrale s'énonce ainsi : comment déployer un modèle
d'apprentissage automatique, dont la taille peut atteindre plusieurs centaines
de mégaoctets, sur un microcontrôleur disposant seulement de quelques centaines
de kilooctets de mémoire vive ? La réponse repose sur trois piliers : un
pipeline de développement rigoureusement structuré, l'application de techniques
avancées de compression des modèles, et une exécution optimisée sur le matériel
cible.

% ---------------------------------------------------------------------
\section{Pipeline de développement TinyML}

Le TinyML ne consiste pas en l'entraînement direct des modèles sur le
microcontrôleur. Le dispositif embarqué n'effectue qu'une seule opération :
l'\textbf{inférence}, c'est-à-dire l'exécution d'un modèle préalablement
entraîné et optimisé. L'intégralité des calculs intensifs s'effectue en amont,
sur des stations de travail ou dans le cloud~\cite{warden2019tinyml}.

\begin{figure}[H]
\centering
\begin{tikzpicture}[node distance=0.45cm]
  \node[bloc, minimum width=2.1cm] (a) {Collecte\\\footnotesize(données)};
  \node[bloc, minimum width=2.1cm, right=of a] (b) {Prétraitement\\\footnotesize(nettoyage)};
  \node[bloco, minimum width=2.1cm, right=of b] (c) {Entraînement\\\footnotesize(PC / cloud)};
  \node[bloco, minimum width=2.1cm, right=of c] (d) {Optimisation\\\footnotesize(compression)};
  \node[blocv, minimum width=2.1cm, below=1.1cm of d] (e) {Déploiement\\\footnotesize(\code{.tflite} $\to$ C)};
  \node[blocv, minimum width=2.1cm, left=of e] (f) {Inférence\\\footnotesize(MCU)};
  \node[blocv, minimum width=2.1cm, left=of f] (g) {Action\\\footnotesize(résultat)};
  \draw[fleche] (a) -- (b);
  \draw[fleche] (b) -- (c);
  \draw[fleche] (c) -- (d);
  \draw[fleche] (d) -- (e);
  \draw[fleche] (e) -- (f);
  \draw[fleche] (f) -- (g);
  \begin{scope}[on background layer]
    \node[draw=orange, dashed, rounded corners=6pt, fit=(a)(d), inner sep=6pt,
          label={[font=\footnotesize\itshape,text=orange]above:sur ordinateur ou dans le cloud}] {};
    \node[draw=vert, dashed, rounded corners=6pt, fit=(g)(e), inner sep=6pt,
          label={[font=\footnotesize\itshape,text=vert]below:sur le microcontrôleur}] {};
  \end{scope}
\end{tikzpicture}
\caption{Pipeline complet de développement TinyML.}
\label{fig:pipeline}
\end{figure}

Le processus se décompose en cinq étapes séquentielles :
\begin{enumerate}
  \item \textbf{Collecte et préparation des données} : l'acquisition de données
        représentatives constitue l'étape fondatrice. Les données proviennent
        généralement des capteurs du dispositif final. Cette phase représente
        une part majeure du temps total de développement, car la qualité d'un
        modèle est directement dépendante de la qualité des données utilisées.
  \item \textbf{Entraînement du modèle} : cette phase s'effectue sur des
        machines dotées de ressources computationnelles importantes, utilisant
        des frameworks tels que TensorFlow ou PyTorch. Le modèle produit est
        généralement trop volumineux (plusieurs mégaoctets en Float32) pour les
        contraintes des microcontrôleurs.
  \item \textbf{Optimisation du modèle} : étape spécifique et critique du
        TinyML, elle applique des techniques de compression (quantification,
        élagage, distillation) pour réduire la taille tout en préservant les
        performances. Le modèle est converti dans le format FlatBuffer de
        TensorFlow Lite (\code{.tflite}).
  \item \textbf{Intégration dans le firmware} : le modèle optimisé est intégré
        au code embarqué sous forme de tableau d'octets statique stocké dans la
        mémoire flash.
  \item \textbf{Exécution de l'inférence} : le microcontrôleur exécute le modèle
        sur les données acquises par ses capteurs, sans connexion réseau. Les
        résultats déclenchent des actions locales.
\end{enumerate}

\textbf{Caractéristique fondamentale :} le modèle est figé au moment de la
compilation. Le microcontrôleur n'effectue aucune forme d'apprentissage ; il
applique les calculs définis.

% ---------------------------------------------------------------------
\section{Techniques de compression et d'optimisation des modèles}

\subsection{Introduction à la compression}

La compression des modèles représente un domaine essentiel pour le déploiement
de l'IA sur dispositifs embarqués~\cite{sze2020efficient}. L'objectif consiste
à réduire significativement la taille des modèles tout en minimisant la
dégradation des performances. Cette réduction permet de respecter les
contraintes mémorielles, d'améliorer l'efficacité énergétique et la vitesse
d'exécution.

Les techniques peuvent être classées en plusieurs catégories : quantification,
réduction de complexité structurelle (élagage, distillation), factorisation et
approches hybrides~\cite{han2016deep}.

\subsection{La quantification (\emph{quantization})}

\subsubsection{Définition et principes}

La quantification est la technique de compression la plus répandue et la plus
efficace en TinyML~\cite{jacob2018quantization,gholami2021survey}. Elle réduit
la précision des représentations numériques des poids et des activations du
réseau.

\textbf{Contexte théorique :} en apprentissage classique, les réseaux
manipulent des valeurs en virgule flottante 32~bits (Float32), offrant une
grande précision mais au prix d'un espace mémoire important (4~octets par
nombre). Dans l'embarqué, cette précision est généralement superflue.

\textbf{Principe :} la quantification consiste à représenter les valeurs sur
8~bits entiers (Int8), divisant par quatre la taille mémoire requise. Un modèle
de 4~Mo en Float32 n'occupera plus que 1~Mo en Int8.

\subsubsection{Mécanismes}

La quantification implique un processus de mise en correspondance
(\emph{mapping}) préservant l'information. Avant d'expliciter ce mécanisme, il
est nécessaire de définir ce qu'est un \textbf{tenseur}.

\paragraph{Qu'est-ce qu'un tenseur ?} En apprentissage automatique, un tenseur
est une généralisation des matrices et des vecteurs à un nombre quelconque de
dimensions. C'est la structure de données fondamentale utilisée pour
représenter les poids, les activations et les données d'entrée dans les réseaux
de neurones. Un tenseur d'ordre $n$ est un tableau à $n$ indices ; chaque
élément est identifié par $n$ coordonnées $(i_1, i_2, \ldots, i_n)$.

\begin{itemize}
  \item \textbf{Ordre 0} (scalaire) : une simple valeur numérique, par exemple $3{,}14$ ;
  \item \textbf{Ordre 1} (vecteur) : un tableau 1D, par exemple les biais d'une couche : $[0{,}2;\ -0{,}5;\ 1{,}3;\ 0{,}8]$ ;
  \item \textbf{Ordre 2} (matrice) : un tableau 2D, comme les poids d'une couche entièrement connectée ;
  \item \textbf{Ordre 3} : un cube de données, typiquement les activations d'une couche convolutionnelle avec plusieurs canaux ;
  \item \textbf{Ordre 4} : les poids des couches convolutionnelles (hauteur $\times$ largeur $\times$ canaux d'entrée $\times$ canaux de sortie).
\end{itemize}

Dans le contexte du TinyML, les \textbf{poids} du réseau sont stockés sous forme
de tenseurs (généralement d'ordre 2 ou 4), les \textbf{activations} (sorties
intermédiaires de chaque couche) sont également des tenseurs, et la
\textbf{quantification} s'applique à chaque tenseur individuellement, avec ses
propres paramètres (\emph{scale} et \emph{zero\_point}).

\paragraph{Formulation mathématique.} La relation fondamentale de la
quantification affine s'exprime par~\cite{jacob2018quantization} :
\begin{equation}
  r \approx S \times (q - Z)
  \label{eq:quant}
\end{equation}
où $r$ est la valeur réelle (Float32), $q$ la valeur quantifiée (entier Int8
dans $[-128, 127]$), $S$ le facteur d'échelle (\emph{scale}, réel positif) et
$Z$ le point zéro (\emph{zero\_point}, entier de même type que $q$).

\paragraph{Processus détaillé.}
\begin{enumerate}
  \item \textbf{Analyse de la plage dynamique} : pour chaque tenseur, on
        détermine les valeurs minimale et maximale ;
  \item \textbf{Calcul des paramètres} : on calcule $S$ et $Z$ pour faire
        correspondre la plage du tenseur à $[-128, 127]$ ;
  \item \textbf{Application de la quantification} : chaque valeur est convertie
        selon la formule~\eqref{eq:quant} inversée ;
  \item \textbf{Exécution en arithmétique entière} : pendant l'inférence, les
        multiplications se font en Int8 et les accumulations en Int32.
\end{enumerate}

\paragraph{Exemple détaillé.} Considérons un tenseur de poids d'une couche
dense avec les valeurs suivantes (en Float32) :
\[
  W = [0{,}23;\ -0{,}45;\ 0{,}12;\ -0{,}87;\ 0{,}65;\ -0{,}33]
\]

\textit{Étape 1 -- Détermination de la plage :} $\min = -0{,}87$,
$\max = 0{,}65$, d'où $\text{plage} = 0{,}65 - (-0{,}87) = 1{,}52$.

\textit{Étape 2 -- Calcul du facteur d'échelle :} sur 8~bits, l'intervalle
$[-128, 127]$ compte 255 pas :
\[
  S = \frac{\text{plage}}{255} = \frac{1{,}52}{255} \approx 0{,}0059608
\]

\textit{Étape 3 -- Calcul du point zéro :} on choisit $Z$ de sorte que la
valeur minimale corresponde à $-128$ :
\[
  Z = -128 - \operatorname{round}\!\left(\frac{\min}{S}\right)
    = -128 - \operatorname{round}(-145{,}95) = -128 + 146 = 18
\]

\textit{Étape 4 -- Quantification de chaque valeur :}
\[
  q_i = \operatorname{round}\!\left(\frac{w_i}{S}\right) + Z
  \quad\Rightarrow\quad
  W_q = [57;\ -57;\ 38;\ -128;\ 127;\ -37]
\]
Toutes les valeurs sont bien comprises dans l'intervalle Int8.

\textit{Étape 5 -- Vérification par déquantification :} pour la première valeur,
\[
  r \approx 0{,}0059608 \times (57 - 18) = 0{,}0059608 \times 39 \approx 0{,}2325
\]
soit une erreur de $0{,}0025$ par rapport à la valeur d'origine $0{,}23$.

\begin{table}[H]
\centering
\small
\begin{tabular}{rrrr}
\toprule
Valeur d'origine & Valeur quantifiée & Valeur déquantifiée & Erreur absolue \\
\midrule
$0{,}23$  & $57$   & $0{,}2325$  & $0{,}0025$ \\
$-0{,}45$ & $-57$  & $-0{,}4471$ & $0{,}0029$ \\
$0{,}12$  & $38$   & $0{,}1192$  & $0{,}0008$ \\
$-0{,}87$ & $-128$ & $-0{,}8703$ & $0{,}0003$ \\
$0{,}65$  & $127$  & $0{,}6497$  & $0{,}0003$ \\
$-0{,}33$ & $-37$  & $-0{,}3278$ & $0{,}0022$ \\
\bottomrule
\end{tabular}
\caption{Exemple de quantification Int8 d'un tenseur de six poids ($S \approx 0{,}0059608$, $Z = 18$).}
\label{tab:quant-exemple}
\end{table}

\textbf{Conclusion de l'exemple :} l'erreur d'arrondi est toujours inférieure
à $S/2 \approx 0{,}003$. Elle est généralement acceptable pour les réseaux de
neurones : la perte de précision globale du modèle reste le plus souvent de
l'ordre de 1 à 2~points~\cite{jacob2018quantization,krishnamoorthi2018}.

\subsubsection{Avantages}
\begin{itemize}
  \item \textbf{Réduction de taille} : division par quatre de la taille des
        poids. Pour un modèle de 20~Mo, cela représente une économie de 15~Mo.
  \item \textbf{Amélioration des performances} : de nombreux microcontrôleurs
        exécutent les opérations entières bien plus rapidement que les opérations
        flottantes. Les opérations Int8 sont particulièrement efficaces sur les
        architectures disposant d'instructions SIMD.
  \item \textbf{Réduction de la consommation énergétique} : une addition 8~bits
        entière consomme environ 30 fois moins d'énergie qu'une addition
        flottante 32~bits, et une multiplication environ 18 fois
        moins~\cite{horowitz2014energy}.
  \item \textbf{Compatibilité avec le matériel} : de nombreux accélérateurs d'IA
        (NPU) sont optimisés pour les calculs Int8.
\end{itemize}

\subsubsection{Méthodes de quantification}

\textbf{Quantification post-entraînement (PTQ)} : convertit un modèle déjà
entraîné. Rapide et simple, elle suffit dans la majorité des cas. La perte de
précision est généralement inférieure à 1--2~\% pour la
classification~\cite{krishnamoorthi2018}.

\textbf{Quantification consciente de l'entraînement (QAT)} : intègre la
quantification dès la phase d'entraînement via une \emph{fake quantization}.
Elle permet d'atteindre des quantifications plus agressives (Int4, binaire)
avec une meilleure préservation de la précision~\cite{jacob2018quantization}.

\begin{lstlisting}[language=Python, caption={Quantification post-entraînement Int8 avec TensorFlow Lite.}, label={lst:ptq}]
import tensorflow as tf

converter = tf.lite.TFLiteConverter.from_saved_model(saved_model_dir)
converter.optimizations = [tf.lite.Optimize.DEFAULT]
# Jeu de données représentatif pour calibrer scale et zero_point
converter.representative_dataset = generateur_donnees_representatives
converter.target_spec.supported_ops = [tf.lite.OpsSet.TFLITE_BUILTINS_INT8]
converter.inference_input_type = tf.int8
converter.inference_output_type = tf.int8
quantized_model = converter.convert()
open("modele_int8.tflite", "wb").write(quantized_model)
\end{lstlisting}

\subsubsection{Limites et défis}
\begin{itemize}
  \item \textbf{Perte de précision} : généralement faible pour la
        classification, elle peut être significative pour la détection d'objets
        ou les opérations sensibles.
  \item \textbf{Gestion des valeurs aberrantes} (\emph{outliers}) : les valeurs
        extrêmes élargissent la plage et dégradent la résolution des autres.
  \item \textbf{Solutions} : quantification par canal, écrêtage
        (\emph{clipping}), utilisation de la QAT.
  \item \textbf{Compatibilité} : certains opérateurs ne sont pas disponibles en
        Int8 dans les moteurs d'inférence embarqués.
\end{itemize}

\subsubsection{Ordres de grandeur}

\begin{table}[H]
\centering
\small
\begin{tabular}{lcccc}
\toprule
Modèle & Taille Float32 & Taille Int8 & Perte de précision & Gain de vitesse \\
\midrule
MobileNetV1  & 16,9 Mo & 4,3 Mo  & $-0{,}8$ \% & $\times 3{,}2$ \\
MobileNetV2  & 14,0 Mo & 3,5 Mo  & $-1{,}2$ \% & $\times 4{,}1$ \\
InceptionV3  & 92,0 Mo & 23,0 Mo & $-1{,}5$ \% & $\times 2{,}8$ \\
\bottomrule
\end{tabular}
\caption{Comparaison indicative des performances avant et après quantification Int8 (ordres de grandeur ; voir~\cite{krishnamoorthi2018,jacob2018quantization}).}
\label{tab:bench-quant}
\end{table}

\subsection{L'élagage (\emph{pruning})}

\subsubsection{Définition et motivation}

L'élagage est une technique de compression structurelle qui supprime les
connexions non essentielles~\cite{han2016deep}. Les réseaux entraînés
contiennent de nombreuses connexions dont les poids sont proches de zéro et qui
contribuent peu à la décision finale.

\textbf{Origine biologique :} l'élagage s'inspire des mécanismes de plasticité
synaptique, dans lesquels les connexions les moins utilisées sont éliminées.

\subsubsection{Méthodologies}

\textbf{Élagage non structuré :} suppression des poids individuellement en
fonction de leur magnitude. Le taux de compression peut être très élevé
(90--95~\%), mais il nécessite des formats de stockage creux spécifiques pour
se traduire par un gain réel.

\textbf{Élagage structuré :} suppression de blocs entiers (filtres, neurones,
canaux). Il réduit effectivement le nombre d'opérations et reste compatible
avec le matériel existant.

\textbf{Algorithme typique :} (1) entraînement initial ; (2) tri des poids par
magnitude ; (3) sélection des $N$~\% de poids les plus petits ; (4) mise à
zéro ; (5) réentraînement court (\emph{fine-tuning}).

\begin{lstlisting}[language=Python, caption={Élagage fondé sur un critère d'importance (approximation de Taylor).}, label={lst:pruning}]
import numpy as np

def taylor_pruning(weights, gradients, pruning_rate=0.5):
    importance = np.abs(weights * gradients)
    threshold = np.percentile(importance, pruning_rate * 100)
    mask = importance > threshold
    return weights * mask, mask
\end{lstlisting}

\subsubsection{Processus complet}

Le processus enchaîne pré-entraînement, élagage (structuré ou non),
\emph{fine-tuning} et itération progressive. \textbf{Stratégie
recommandée :} plusieurs itérations avec des taux modérés (par exemple 20~\% à
chaque itération) plutôt qu'une seule itération avec un taux élevé.

\subsubsection{Résultats expérimentaux}

\begin{table}[H]
\centering
\small
\begin{tabular}{cccc}
\toprule
Taux d'élagage & Précision & Perte & Réduction des paramètres \\
\midrule
20 \% & 70,9 \% & $-0{,}3$ \% & 20 \% \\
40 \% & 70,1 \% & $-1{,}1$ \% & 40 \% \\
60 \% & 68,5 \% & $-2{,}7$ \% & 60 \% \\
80 \% & 65,2 \% & $-6{,}0$ \% & 80 \% \\
\bottomrule
\end{tabular}
\caption{Impact indicatif de l'élagage sur la précision de MobileNetV1 (ImageNet).}
\label{tab:pruning}
\end{table}

\subsection{La distillation (\emph{knowledge distillation})}

\subsubsection{Concept}

La distillation transfère les connaissances d'un modèle complexe (le
\emph{professeur}) vers un modèle plus petit (l'\emph{élève}), permettant de
créer des modèles compacts tout en conservant une grande partie des
performances~\cite{hinton2015distilling}.

\textbf{Analogie pédagogique :} un expert transmet non seulement les réponses
finales, mais également son raisonnement à un élève.

\subsubsection{Principe théorique}

Le transfert de connaissances peut porter sur les sorties (\emph{output
knowledge}), sur les représentations intermédiaires (\emph{intermediate
knowledge}) ou sur les relations entre exemples (\emph{relation knowledge}).
La fonction de perte combine deux termes :
\begin{equation}
  L_{\text{total}} = \alpha \, L_{\text{classification}} + \beta \, L_{\text{distillation}}
\end{equation}
Une \textbf{température} $T > 1$ est utilisée pour adoucir les distributions de
probabilités produites par le professeur~\cite{hinton2015distilling} :
\begin{equation}
  P_i = \frac{\exp(z_i / T)}{\sum_j \exp(z_j / T)}
\end{equation}

\subsubsection{Méthodes}
\begin{itemize}
  \item \textbf{Distillation standard} : l'élève apprend à reproduire les sorties du professeur.
  \item \textbf{Distillation par représentations intermédiaires} : l'élève apprend les représentations de plusieurs couches internes.
  \item \textbf{Distillation mutuelle} : plusieurs élèves s'enseignent mutuellement.
  \item \textbf{Auto-distillation} : un modèle enseigne à une version plus petite de lui-même.
\end{itemize}

\subsubsection{Résultats}

\begin{table}[H]
\centering
\small
\begin{tabular}{llcccc}
\toprule
Professeur & Élève & Param. P & Param. E & Précision P & Précision E \\
\midrule
ResNet-152 & ResNet-18 & 60,2 M & 11,7 M & 78,3 \% & 74,8 \% \\
ResNet-152 & MobileNet & 60,2 M & 4,2 M  & 78,3 \% & 71,3 \% \\
\bottomrule
\end{tabular}
\caption{Comparaison indicative des performances (top-1 ImageNet) avant et après distillation ; M = millions de paramètres.}
\label{tab:distillation}
\end{table}

\subsection{Combinaison des techniques}

Les techniques sont complémentaires et sont généralement appliquées dans
l'ordre suivant~\cite{han2016deep} : (1) distillation pour réduire la taille de
l'architecture ; (2) élagage pour réduire la complexité ; (3) quantification
pour réduire la précision numérique ; avec un \emph{fine-tuning} après chaque
étape.

\begin{figure}[H]
\centering
\begin{tikzpicture}[node distance=0.5cm]
  \node[bloco, minimum width=2.3cm] (m) {Modèle Float32\\\footnotesize(plusieurs Mo)};
  \node[bloc, minimum width=2.1cm, right=of m] (d) {Distillation};
  \node[bloc, minimum width=2.1cm, right=of d] (e) {Élagage\\+ \emph{fine-tuning}};
  \node[bloc, minimum width=2.1cm, right=of e] (q) {Quantification\\Int8};
  \node[blocv, minimum width=2.3cm, right=of q] (f) {Modèle TinyML\\\footnotesize(quelques dizaines de Ko)};
  \draw[fleche] (m) -- (d);
  \draw[fleche] (d) -- (e);
  \draw[fleche] (e) -- (q);
  \draw[fleche] (q) -- (f);
\end{tikzpicture}
\caption{Enchaînement typique des techniques de compression.}
\label{fig:compression}
\end{figure}

\begin{table}[H]
\centering
\small
\begin{tabular}{lccc}
\toprule
Méthode & Réduction de taille & Perte de précision & Gain de vitesse \\
\midrule
Quantification seule & $\times 4$         & $< 1$ \%     & $\times 2$ à $4$ \\
Élagage seul         & $\times 5$ à $10$  & 2 à 5 \%     & $\times 1$ à $2$ \\
Distillation seule   & $\times 10$ à $50$ & 3 à 8 \%     & $\times 1$ à $3$ \\
Combinaison          & $\times 100$ à $300$ & 5 à 15 \%  & $\times 5$ à $15$ \\
\bottomrule
\end{tabular}
\caption{Efficacité comparée des méthodes de compression (ordres de grandeur).}
\label{tab:compression}
\end{table}

\subsection{Techniques avancées}
\begin{itemize}
  \item \textbf{Factorisation de matrices} : décomposition des matrices de poids
        en produits de matrices de plus petites dimensions (SVD, rang faible).
  \item \textbf{Binarisation} : représentation des poids et activations par des
        valeurs binaires ($-1$ ou $+1$). La réduction de taille est extrême
        ($\times 32$) mais la perte de précision importante.
  \item \textbf{Réseaux creux} : utilisation de représentations creuses pour les
        matrices de poids (formats COO, CSR).
  \item \textbf{Recherche automatique d'architecture (NAS)} : conception
        automatique de réseaux adaptés à l'embarqué, comme MobileNet,
        EfficientNet ou MCUNet~\cite{lin2020mcunet}.
\end{itemize}

\subsection{Métriques d'évaluation}
\begin{itemize}
  \item \textbf{Taux de compression} : $\text{Taille}_{\text{originale}} / \text{Taille}_{\text{compressée}}$.
  \item \textbf{Métriques de performance} : précision, latence, consommation énergétique, mémoire RAM et flash.
  \item \textbf{Bancs d'essai} : ImageNet, CIFAR-10/100 et, pour l'embarqué, MLPerf Tiny~\cite{banbury2021mlperftiny}.
\end{itemize}

\subsection{Études de cas}
\begin{itemize}
  \item \textbf{Détection de mots-clés} : un modèle CNN de l'ordre du mégaoctet
        peut être ramené à une centaine de kilooctets, avec une précision
        supérieure à 90~\% et une latence de quelques dizaines de millisecondes
        sur microcontrôleur~\cite{banbury2021mlperftiny}.
  \item \textbf{Classification d'images} : des réseaux de type MobileNet ou
        MCUNet, quantifiés en Int8, tiennent dans moins de 1~Mo et s'exécutent
        sur des microcontrôleurs haut de gamme~\cite{lin2020mcunet}.
  \item \textbf{Séries temporelles} (vibrations, capteurs inertiels) : de petits
        réseaux denses ou auto-encodeurs de quelques dizaines de kilooctets
        suffisent pour la détection d'anomalies, avec une latence de quelques
        millisecondes sur ESP32 ; c'est le type de modèle visé par le scénario illustratif de la partie~\ref{ch:fondements}.
\end{itemize}

% =====================================================================
%  PARTIE 3 — Matériel et outils de développement
% =====================================================================
\chapter{Matériel et outils de développement}
\label{ch:materiel}

Les techniques de compression ne suffisent pas à elles seules : leur efficacité
dépend du matériel sur lequel le modèle est exécuté et des outils logiciels qui
assurent la conversion et le déploiement. Cette partie présente trois
architectures représentatives de l'évolution des microcontrôleurs pour le
TinyML, puis la chaîne d'outils logiciels associée.

\section{Architectures matérielles}

\subsection{ESP32-S3 (Espressif)}
\begin{itemize}
  \item \textbf{Cœur} : Xtensa LX7 double cœur à 240~MHz ;
  \item \textbf{Mémoire} : 512~Ko de SRAM, jusqu'à 8~Mo de PSRAM externe ;
  \item \textbf{Accélération} : extensions vectorielles (SIMD) pour les calculs 8/16~bits ;
  \item \textbf{Connectivité} : Wi-Fi et Bluetooth LE intégrés~\cite{espressif-esp32s3}.
\end{itemize}

\subsection{ESP32-P4 (Espressif)}
\begin{itemize}
  \item \textbf{Architecture} : RISC-V double cœur jusqu'à 400~MHz ;
  \item \textbf{Mémoire} : jusqu'à 32~Mo de PSRAM ;
  \item \textbf{Accélération} : extensions dédiées à l'IA, traitement d'image matériel ;
  \item \textbf{Sortie} : 2024~\cite{espressif-esp32p4}.
\end{itemize}

\subsection{STM32N6 (STMicroelectronics)}
\begin{itemize}
  \item \textbf{Cœur} : Arm Cortex-M55 à 800~MHz ;
  \item \textbf{NPU} : accélérateur neuronal Neural-ART à 600~GOPS ;
  \item \textbf{Efficacité} : environ 3~TOPS/W ;
  \item \textbf{Mémoire} : 4,2~Mo de SRAM intégrée~\cite{st-stm32n6}.
\end{itemize}

\begin{table}[H]
\centering
\small
\begin{tabularx}{\linewidth}{>{\bfseries}l X X X}
\toprule
Critère & ESP32-S3 & ESP32-P4 & STM32N6 \\
\midrule
Cœur           & Xtensa LX7, 2 cœurs, 240 MHz & RISC-V, 2 cœurs, 400 MHz & Cortex-M55, 800 MHz \\
SRAM / PSRAM   & 512 Ko / jusqu'à 8 Mo & 768 Ko / jusqu'à 32 Mo & 4,2 Mo SRAM \\
Accélération IA & Instructions SIMD & Extensions IA, traitement d'image & NPU Neural-ART (600 GOPS) \\
Connectivité   & Wi-Fi, BLE & Aucune radio intégrée & Aucune radio intégrée \\
Usage typique  & Audio, capteurs, petite vision & Vision, interface homme-machine & Vision temps réel, modèles plus lourds \\
\bottomrule
\end{tabularx}
\caption{Comparaison des architectures matérielles pour le TinyML.}
\label{tab:materiel}
\end{table}

Ces trois architectures illustrent l'évolution des microcontrôleurs pour le
TinyML, avec une complexité et des performances croissantes : des instructions
vectorielles d'un MCU généraliste jusqu'à l'intégration d'un véritable
accélérateur neuronal. Pour un scénario comme la détection d'anomalies
vibratoires, un ESP32 ou un ESP32-S3 est largement suffisant ; le Wi-Fi
intégré permet en outre de transmettre les alertes.

\section{Chaîne d'outils logiciels}

\begin{figure}[H]
\centering
\begin{tikzpicture}[node distance=0.6cm]
  \node[bloco, minimum width=2.8cm] (a) {1. Entraînement\\\footnotesize(TensorFlow / PyTorch)};
  \node[bloc, minimum width=2.8cm, right=of a] (b) {2. Conversion\\\footnotesize(TFLite Converter)};
  \node[blocv, minimum width=2.8cm, right=of b] (c) {3. Exécution\\\footnotesize(TFLM / ESP-DL)};
  \node[blocv, minimum width=2.8cm, right=of c] (d) {4. Firmware\\\footnotesize(ESP-IDF / CubeIDE)};
  \draw[fleche] (a) -- (b);
  \draw[fleche] (b) -- (c);
  \draw[fleche] (c) -- (d);
\end{tikzpicture}
\caption{Chaîne d'outils logiciels pour le développement TinyML.}
\label{fig:outils}
\end{figure}

\subsection{TensorFlow Lite for Microcontrollers (TFLM)}

TFLM est le moteur d'inférence de référence pour le TinyML~\cite{david2021tflm,tflm-doc} :
\begin{itemize}
  \item interpréteur C++ minimaliste, fonctionnant sans système d'exploitation
        ni allocation dynamique de mémoire ;
  \item empreinte de quelques dizaines de kilooctets seulement ;
  \item modèle intégré sous forme de tableau C dans le firmware ;
  \item \emph{tensor arena} : zone de mémoire statique, réservée à la
        compilation, qui contient les activations pendant l'inférence.
\end{itemize}

La conversion d'un fichier \code{.tflite} en tableau C se fait en une ligne de
commande, puis le modèle est chargé par l'interpréteur (listing~\ref{lst:tflm}).

\begin{lstlisting}[language=C++, caption={Initialisation de l'interpréteur TFLM sur microcontrôleur.}, label={lst:tflm}]
// Généré par : xxd -i modele_int8.tflite > modele_data.cc
#include "modele_data.h"
#include "tensorflow/lite/micro/micro_interpreter.h"
#include "tensorflow/lite/micro/micro_mutable_op_resolver.h"

constexpr int kTensorArenaSize = 20 * 1024;      // 20 Ko de SRAM
alignas(16) static uint8_t tensor_arena[kTensorArenaSize];

const tflite::Model* model = tflite::GetModel(modele_data);
static tflite::MicroMutableOpResolver<3> resolver;
resolver.AddFullyConnected();
resolver.AddRelu();
resolver.AddSoftmax();

static tflite::MicroInterpreter interpreter(model, resolver,
                                            tensor_arena, kTensorArenaSize);
interpreter.AllocateTensors();
\end{lstlisting}

\subsection{Edge Impulse}

Edge Impulse est une plateforme en ligne couvrant l'ensemble du pipeline, de la
collecte des données au déploiement~\cite{edgeimpulse}. Elle permet
d'enregistrer des données depuis la carte, de concevoir le prétraitement
(FFT, MFCC) et le modèle, d'évaluer ses performances estimées sur le
microcontrôleur cible, puis de générer automatiquement une bibliothèque C++ ou
Arduino prête à être compilée.

\subsection{Outils spécifiques aux fabricants}
\begin{itemize}
  \item \textbf{ESP-DL} (Espressif) : bibliothèque optimisée pour les ESP32, qui
        exploite leurs instructions vectorielles~\cite{espressif-espdl} ;
  \item \textbf{STM32Cube.AI} (STMicroelectronics) : outil de conversion et
        d'optimisation des modèles pour les STM32~\cite{st-cubeai} ;
  \item \textbf{NXP eIQ} : environnement de développement pour les
        microcontrôleurs et processeurs NXP~\cite{nxp-eiq}.
\end{itemize}

% =====================================================================
%  PARTIE 4 — Bénéfices et domaines d'application
% =====================================================================
\chapter{Bénéfices et domaines d'application}
\label{ch:benefices}

\section{Les bénéfices de l'AIoT et du TinyML}

L'association de l'intelligence artificielle et de l'Internet des objets,
appelée AIoT, permet aux équipements connectés d'analyser les données qu'ils
collectent afin de reconnaître des situations ou de détecter des anomalies. Le
TinyML correspond plus précisément à l'exécution de modèles d'apprentissage
automatique sur des dispositifs disposant de ressources limitées, notamment des
microcontrôleurs. Cette approche déplace une partie de l'intelligence au plus
près des capteurs, plutôt que de transmettre systématiquement les données à un
serveur distant. Elle présente plusieurs avantages en matière de réactivité, de
confidentialité et de consommation énergétique~\cite{ray2022review}.

\subsection{Une meilleure réactivité et une moindre dépendance au réseau}

Dans une architecture reposant sur le cloud, les données doivent être
transmises à un serveur avant d'être analysées. Le délai de réponse dépend donc
non seulement du calcul, mais aussi de la connexion réseau et des échanges avec
le serveur. Avec le TinyML, l'analyse peut être réalisée directement dans
l'objet connecté, ce qui évite cet aller-retour pour les décisions prises
localement.

Cette organisation est particulièrement intéressante lorsqu'un système doit
réagir rapidement, par exemple pour reconnaître un événement sonore ou détecter
un comportement inhabituel d'une machine. Elle permet également de maintenir
certaines fonctions d'analyse lorsque la connexion Internet est indisponible, à
condition que ces fonctions aient été conçues pour fonctionner de manière
autonome.

Il serait toutefois incorrect de parler de \og latence zéro\fg{}.
L'acquisition des données, leur préparation et l'exécution du modèle
nécessitent toujours du temps. Le bénéfice correspond à une réduction de la
latence et de sa dépendance au réseau, et non à sa disparition.

\subsection{Une meilleure maîtrise de la confidentialité des données}

Les objets connectés peuvent collecter des informations sensibles, comme des
enregistrements audio, des données de mouvement ou des mesures physiologiques.
Lorsque ces données sont systématiquement envoyées vers un serveur, leur
traitement implique davantage de transferts et de stockage à distance.
L'inférence locale permet de limiter ces échanges en analysant les informations
directement sur l'appareil.

Par exemple, un dispositif de reconnaissance sonore peut analyser un signal
audio localement et ne transmettre que le résultat de la détection. Selon
l'architecture retenue, il devient ainsi possible de communiquer une alerte
plutôt que l'ensemble de l'enregistrement. Le TinyML contribue donc à réduire
la quantité de données brutes qui quittent l'objet connecté.

Cet avantage ne signifie pas que la confidentialité est automatiquement
garantie. Certains systèmes peuvent encore transmettre des données pour
l'entraînement, le suivi ou d'autres services. Il faut donc distinguer la
possibilité de traiter les données localement du choix réel de ne pas les
transmettre.

\subsection{Une réduction possible de la consommation énergétique}

L'autonomie constitue une contrainte importante pour les objets connectés
alimentés par batterie. Les communications sans fil consomment de l'énergie,
notamment lorsque l'appareil transmet fréquemment des mesures. En analysant les
données localement et en ne communiquant que les résultats utiles, le TinyML
peut réduire le recours à la transmission radio~\cite{warden2019tinyml}.

Cependant, le calcul local consomme lui aussi de l'énergie. Le gain final
dépend donc de plusieurs paramètres : la complexité du modèle, la fréquence des
analyses, la consommation des capteurs, les modes de veille et la fréquence des
communications. Une analyse locale bien dimensionnée peut améliorer
l'autonomie, mais cet avantage doit être mesuré sur le dispositif complet
plutôt que supposé à partir du seul modèle d'IA~\cite{banbury2021mlperftiny}.

\begin{table}[H]
\centering
\small
\begin{tabularx}{\linewidth}{>{\bfseries}l X X}
\toprule
Bénéfice & Ce que le TinyML apporte & Nuance à garder en tête \\
\midrule
Réactivité      & Décision locale, sans aller-retour réseau & La latence est réduite, pas supprimée \\
Fonctionnement hors ligne & Analyse maintenue sans connexion & Seulement pour les fonctions conçues pour cela \\
Confidentialité & Les données brutes peuvent rester dans l'objet & Dépend des choix réels de transmission \\
Énergie         & Moins de transmissions radio & Le calcul local consomme aussi \\
Coût            & Moins de données envoyées et stockées dans le cloud & Coût de développement et de mise à jour \\
\bottomrule
\end{tabularx}
\caption{Synthèse des bénéfices du TinyML et de leurs limites.}
\label{tab:benefices}
\end{table}

\section{Domaines d'application et exemples d'utilisation}

\subsection{Industrie : surveillance des machines et maintenance}

Dans l'industrie, l'apprentissage automatique embarqué peut servir à surveiller
le fonctionnement des machines et à détecter des signes de dégradation. La
maintenance prédictive fait notamment partie des applications de l'IA en
périphérie : l'objectif est d'identifier des défauts avant qu'ils n'entraînent
des problèmes plus importants~\cite{zhang2021aiot}.

Le scénario illustratif présenté dans la partie~\ref{ch:fondements} relève de
ce domaine : un moteur surveillé par un capteur de vibrations, dont les mesures
seraient analysées sur un microcontrôleur afin de déclencher une alerte
lorsqu'un comportement inhabituel est détecté.

L'intérêt recherché est de fournir une information utile au personnel de
maintenance sans transmettre continuellement toutes les mesures à un serveur.
Il convient néanmoins de distinguer la détection d'anomalies de la prédiction
d'une panne : repérer un fonctionnement inhabituel ne permet pas nécessairement
de déterminer sa cause ou le temps restant avant une défaillance. La détection
constitue un premier niveau de surveillance ; la prédiction demande une analyse
et une validation supplémentaires.

\subsection{Maison connectée : reconnaissance de mots-clés et de sons}

La reconnaissance de mots-clés est une application représentative du TinyML ;
elle fait d'ailleurs partie des tâches de référence du banc d'essai MLPerf
Tiny~\cite{banbury2021mlperftiny}. Un modèle compact peut être chargé
d'identifier un nombre limité de commandes vocales sur un appareil disposant de
peu de mémoire et de puissance de calcul. Cette tâche ciblée se distingue d'une
conversation complète avec un assistant fondé sur un grand modèle de langage.

L'analyse locale présente ici un double intérêt : réduire les échanges avec un
serveur et permettre une réaction rapide. Elle peut aussi limiter la
transmission de données audio lorsque le système est conçu pour ne communiquer
que les résultats nécessaires.

La reconnaissance reste cependant sensible aux conditions d'utilisation. Le
bruit ambiant peut perturber la détection et provoquer des erreurs. Pour
obtenir un système robuste, il faut donc entraîner et tester le modèle avec des
données représentatives des environnements dans lesquels l'objet sera
réellement utilisé.

% =====================================================================
%  PARTIE 5 — Limites, défis et perspectives
% =====================================================================
\chapter{Limites, défis et perspectives d'avenir}
\label{ch:limites}

\section{Des ressources matérielles limitées}

La principale difficulté du TinyML consiste à faire fonctionner un modèle utile
sur un appareil disposant de ressources réduites. La mémoire vive, l'espace de
stockage, la puissance de calcul et l'énergie disponible limitent la complexité
des modèles déployables. L'optimisation doit ainsi prendre en compte plusieurs
critères simultanément : la taille du modèle, sa précision, son temps
d'exécution et sa consommation~\cite{banbury2020benchmarking}.

Les outils logiciels imposent également des contraintes. TensorFlow Lite for
Microcontrollers, par exemple, prend en charge un ensemble limité d'opérations
et nécessite une gestion adaptée de la mémoire~\cite{david2021tflm}. Un modèle
conçu sur ordinateur ne peut donc pas toujours être transféré directement sur un
microcontrôleur sans modification.

Ces contraintes conduisent à rechercher un compromis : le modèle le plus
complexe n'est pas nécessairement le plus pertinent si un modèle plus simple
répond au besoin tout en consommant moins de ressources.

\section{Une fiabilité dépendante des données et du contexte}

Les performances d'un modèle dépendent des données utilisées pour son
entraînement et de leur représentativité. Dans les applications vocales, le
bruit peut dégrader la reconnaissance. Dans les dispositifs portables, les
mouvements, les différences entre utilisateurs ou les variations de
positionnement des capteurs peuvent modifier les signaux analysés.

Un modèle peut donc donner de bons résultats pendant les essais puis devenir
moins fiable lorsque les conditions changent. Il est nécessaire de le tester
dans des situations variées et de suivre ses performances après le
déploiement. Cette exigence est particulièrement importante lorsque les erreurs
peuvent avoir des conséquences sur la surveillance d'un équipement ou sur
l'interprétation de données de santé.

\section{L'entraînement et l'évolution des modèles}

Dans de nombreux systèmes TinyML, l'entraînement est effectué sur un ordinateur
ou une infrastructure distante, tandis que le microcontrôleur réalise uniquement
l'inférence, c'est-à-dire l'utilisation du modèle déjà entraîné. Cette
séparation s'explique par les besoins généralement plus élevés de
l'entraînement en mémoire et en puissance de calcul. TensorFlow Lite for
Microcontrollers ne prend d'ailleurs pas en charge l'entraînement sur
l'appareil~\cite{tflm-doc}.

Cette organisation soulève la question de l'évolution du système. Lorsque les
données ou les conditions d'utilisation changent, une nouvelle version du modèle
peut devenir nécessaire ; elle doit alors être redéployée, par exemple par une
mise à jour du firmware à distance (OTA).

L'apprentissage directement sur l'appareil constitue une piste d'évolution, des
travaux ayant montré qu'il est possible d'entraîner partiellement un réseau
avec moins de 256~Ko de mémoire~\cite{lin2022ondevice}. Il n'est cependant pas
encore une solution générale pour tous les microcontrôleurs : il reste limité
par les ressources disponibles et par la complexité supplémentaire qu'il
introduit.

\section{La sécurité des modèles embarqués}

Placer un modèle dans un objet physique l'expose à de nouvelles menaces : un
attaquant ayant accès à l'appareil peut tenter d'extraire le modèle de la
mémoire flash, de le modifier ou de fournir au capteur des entrées conçues pour
le tromper. La protection de la mémoire (chiffrement de la flash, démarrage
sécurisé) et la sécurisation des mises à jour OTA, que proposent par exemple
les ESP32~\cite{espressif-esp32}, font donc partie intégrante de la conception
d'un système AIoT.

\section{Vers des architectures complémentaires entre l'objet et le cloud}

Le développement du TinyML ne signifie pas nécessairement la disparition du
cloud. Les deux approches peuvent remplir des rôles complémentaires : l'objet
exécute les analyses locales compatibles avec ses ressources, tandis qu'une
infrastructure plus puissante peut prendre en charge l'entraînement et les
traitements plus exigeants.

\begin{figure}[H]
\centering
\begin{tikzpicture}[node distance=2.4cm]
  \node[blocv, text width=2.9cm] (o) {Objets TinyML\\\footnotesize inférence locale, alertes};
  \node[bloc, text width=2.9cm, right=of o] (e) {Passerelle / edge\\\footnotesize agrégation, analyses};
  \node[bloco, text width=2.9cm, right=of e] (c) {Cloud\\\footnotesize entraînement, supervision};
  \draw[fleche] ([yshift=4pt]o.east) -- node[above,font=\scriptsize]{résultats} ([yshift=4pt]e.west);
  \draw[fleche] ([yshift=4pt]e.east) -- node[above,font=\scriptsize]{données utiles} ([yshift=4pt]c.west);
  \draw[fleche, dashed] ([yshift=-4pt]c.west) -- node[below,font=\scriptsize]{nouveau modèle} ([yshift=-4pt]e.east);
  \draw[fleche, dashed] ([yshift=-4pt]e.west) -- node[below,font=\scriptsize]{mise à jour OTA} ([yshift=-4pt]o.east);
\end{tikzpicture}
\caption{Répartition des traitements dans une architecture AIoT hybride.}
\label{fig:hybride}
\end{figure}

Les perspectives reposent notamment sur l'amélioration de l'efficacité des
modèles et de leur adaptation aux contraintes matérielles. Les travaux récents
sur la quantification montrent que la recherche vise à réduire les besoins en
mémoire et en énergie tout en préservant des performances adaptées à
l'application~\cite{gholami2021survey}. Parallèlement, les microcontrôleurs
intègrent de plus en plus d'accélérateurs neuronaux dédiés, comme le montre le
STM32N6 présenté dans la partie~\ref{ch:materiel}.

L'enjeu n'est donc pas de placer toute l'intelligence dans chaque capteur, mais
de répartir les traitements de façon pertinente selon les besoins de
réactivité, d'autonomie et de confidentialité.


% ---------------- Conclusion ----------------
% =====================================================================
%  Conclusion générale
% =====================================================================
\chapter*{Conclusion générale}
\addcontentsline{toc}{chapter}{Conclusion générale}
\markboth{Conclusion générale}{}

Ce rapport a montré comment l'intelligence artificielle peut quitter les
centres de données pour s'installer au cœur des objets connectés. Partant de
l'architecture classique de l'IoT, où l'objet se contente de mesurer et de
transmettre, nous avons vu que l'AIoT donne aux systèmes connectés la capacité
d'interpréter leurs données, et que le TinyML pousse cette logique jusqu'aux
microcontrôleurs les plus modestes.

Cette évolution repose sur une séparation nette entre l'entraînement, réalisé
sur des machines puissantes, et l'inférence, seule exécutée dans l'objet. Elle
n'est possible que grâce à un pipeline de développement rigoureux et à des
techniques de compression (quantification, élagage, distillation) capables de
réduire la taille d'un modèle d'un ou deux ordres de grandeur, ainsi qu'à des
microcontrôleurs et des outils logiciels de plus en plus adaptés.

Les bénéfices sont réels : réactivité accrue, fonctionnement en l'absence de
réseau, réduction des données transmises et donc meilleure maîtrise de la
confidentialité et, souvent, de la consommation énergétique. Ils doivent
cependant être nuancés : la latence n'est pas nulle, le calcul local consomme
lui aussi de l'énergie, la fiabilité dépend fortement des données
d'entraînement et l'évolution des modèles déployés reste délicate.

Le TinyML ne remplace donc pas le cloud : il le complète. L'avenir de l'AIoT se
dessine comme une intelligence distribuée, dans laquelle chaque niveau (objet,
passerelle, cloud) exécute les traitements qui lui conviennent le mieux. Notre
exemple de détection d'anomalies vibratoires sur ESP32 illustre bien cette
répartition : l'objet surveille et alerte en temps réel, tandis que le cloud
supervise, archive et améliore le modèle au fil du temps.


% ---------------- Bibliographie ----------------
\clearpage
\phantomsection
\addcontentsline{toc}{chapter}{Bibliographie et webographie}
\markboth{Bibliographie et webographie}{}
\nocite{*}
\printbibheading[title={Bibliographie et webographie}]
\printbibliography[keyword=ouvrage,heading=subbibliography,title={Ouvrages de référence}]
\printbibliography[keyword=article,heading=subbibliography,title={Articles scientifiques}]
\printbibliography[keyword=rapport,heading=subbibliography,title={Normes, rapports et documentation technique}]
\printbibliography[keyword=web,heading=subbibliography,title={Ressources en ligne et communautés}]

\end{document}
